% Exact LaTeX recreation of the supplied Phase-1.pdf
% Compile with pdfLaTeX. Place member pictures in the same folder if required.

\documentclass[11pt,a4paper]{article}

\usepackage[a4paper,left=2.0cm,right=2.0cm,top=1.55cm,bottom=1.55cm]{geometry}
\usepackage[T1]{fontenc}
\usepackage{lmodern}
\usepackage{graphicx}
\usepackage{xcolor}
\usepackage{array}
\usepackage{tabularx}
\usepackage{ragged2e}
\usepackage{enumitem}
\usepackage{fancyhdr}
\usepackage{tikz}
\usepackage{setspace}

% -------------------------------------------------
% COLOURS
% -------------------------------------------------
\definecolor{TitleBlue}{RGB}{61,103,154}

% -------------------------------------------------
% GENERAL SETTINGS
% -------------------------------------------------
\setlength{\parindent}{0pt}
\setlength{\parskip}{0pt}
\renewcommand{\arraystretch}{1.0}

% The supplied document uses a clean sans-serif heading style
% and italic serif-style explanatory text.
\newcommand{\blueheading}[1]{%
    {\sffamily\color{TitleBlue}\fontsize{15}{18}\selectfont #1}%
}

\newcommand{\bluebigheading}[1]{%
    {\sffamily\bfseries\color{TitleBlue}\fontsize{16}{19}\selectfont #1}%
}

\newcommand{\instruction}[1]{%
    {\itshape\fontsize{10.5}{12.5}\selectfont #1}%
}

\newcommand{\answerbox}[2][3.0cm]{%
    \vspace{2pt}
    \noindent
    \fbox{%
        \begin{minipage}[t][#1][t]{\dimexpr\textwidth-2\fboxsep-2\fboxrule\relax}
            \vspace{1mm}
            \itshape #2
        \end{minipage}
    }
}

% -------------------------------------------------
% HEADER / FOOTER
% -------------------------------------------------
\pagestyle{fancy}
\fancyhf{}
\renewcommand{\headrulewidth}{0pt}
\renewcommand{\footrulewidth}{0pt}

\fancyhead[L]{%
    \sffamily\bfseries\itshape\fontsize{10.5}{12}\selectfont
    DSCPP-III-2026-T265
}
\fancyhead[R]{%
    \sffamily\bfseries\fontsize{10.5}{12}\selectfont
    PROJECT PROPOSAL
}
\fancyfoot[R]{%
    \itshape\fontsize{10}{12}\selectfont
    Page \thepage
}
\setlength{\headheight}{14pt}
\setlength{\headsep}{23pt}
\setlength{\footskip}{24pt}

% -------------------------------------------------
% PAGE 1
% -------------------------------------------------
\begin{document}

\vspace*{4mm}

\bluebigheading{PROJECT AND TEAM INFORMATION}

\vspace{13mm}

\blueheading{Project Title}

\vspace{1mm}

\instruction{(Try to choose a catchy title. Max 20 words).}

\answerbox[1.0cm]{Campus X: Beyond the Feed}

\vspace{13mm}

\blueheading{Student / Team Information}

\vspace{2mm}

% Team information table
% Compact 4-member layout so all members remain on Page 1.
\noindent
\begin{tabularx}{\textwidth}{|>{\RaggedRight\arraybackslash}p{0.69\textwidth}|>{\RaggedRight\arraybackslash}X|}
\hline
\begin{minipage}[t]{\linewidth}
\vspace{1.5mm}
\textit{Team Name:}\\[-1mm]
\textit{Team \#}
\vspace{1.5mm}
\end{minipage}
&
\begin{minipage}[t]{\linewidth}
\vspace{1.5mm}
\textit{Meta Minds}
\vspace{1.5mm}
\end{minipage}
\\
\hline
\begin{minipage}[t][3.25cm][t]{\linewidth}
\vspace{1.5mm}
\textbf{Team member 1 (Team Lead)}\\[-1mm]
\textit{(Last Name, name: student ID: email, picture):}
\vspace{2.0cm}
\end{minipage}
&
\begin{minipage}[t][3.25cm][t]{\linewidth}
\vspace{1.5mm}
\textit{Bhatnagar, Aryaman -- 2510380065}\\
\textit{aryaman6649@gmail.com}

\vspace{2mm}
\includegraphics[width=3.0cm,height=1.5cm,keepaspectratio]{aryaman.png}
\end{minipage}
\\
\hline
\begin{minipage}[t][3.7cm][t]{\linewidth}
\vspace{1.5mm}
\textbf{Team member 2}\\[-1mm]
\textit{(Last Name, name: student ID: email, picture):}
\vspace{2.0cm}
\end{minipage}
&
\begin{minipage}[t][3.7cm][t]{\linewidth}
\vspace{1.5mm}
\textit{Batra, Akshat -- 2510360007}\\
\textit{akshatbatra0616@gmail.com}

\vspace{2mm}
\includegraphics[width=3.2cm,height=1.6cm,keepaspectratio]{akshat.jpg}
\end{minipage}
\\
\hline
\begin{minipage}[t][3.25cm][t]{\linewidth}
\vspace{1.5mm}
\textbf{Team member 3}\\[-1mm]
\textit{(Last Name, name: student ID: email, picture):}
\vspace{2.0cm}
\end{minipage}
&
\begin{minipage}[t][3.7cm][t]{\linewidth}
\vspace{1.5mm}
\textit{Bhadana, Vedant -- 2510370565}\\
\textit{bhadanavedant29
@gmail.com}

\vspace{2mm}
\includegraphics[width=3.0cm,height=1.5cm,keepaspectratio]{vedant.jpeg}
\end{minipage}
\\
\hline
\begin{minipage}[t][3.7cm][t]{\linewidth}
\vspace{1.5mm}
\textbf{Team member 4}\\[-1mm]
\textit{(Last Name, name: student ID: email, picture):}
\vspace{2.0cm}
\end{minipage}
&
\begin{minipage}[t][3.7cm][t]{\linewidth}
\vspace{1.5mm}
\textit{Rastogi, Ritwik -- 2510374220}\\
\textit{ritwik.com2@gmail.com}

\vspace{2mm}
 \includegraphics[width=3.0cm,height=1.5cm,keepaspectratio]{ritwik.jpeg}
\end{minipage}
\\
\hline
\end{tabularx}

\newpage

% -------------------------------------------------
% PAGE 2
% -------------------------------------------------

\bluebigheading{PROPOSAL DESCRIPTION (10 pts)}

\vspace{12mm}

\blueheading{Motivation (1 pt)}

\vspace{1mm}

\instruction{(Describe the problem you want to solve and why it is important. Max 300 words).}

\answerbox[5cm]{We want to build Campus X because we believe college life can be more organized and connected with a platform. Our main goal is to create something that students can actually use every day not a project made for showing off. Campus X will bring together campus activities like announcements, events, clubs, academic resources, student connections, Lost & Found emergency updates and other campus activities all in one place. We want students to find what they need quickly and easily without having to check groups or sources. We also want this project to be practical and meaningful for our learning. Since we are studying Data Structures and OOPs, in C++ Campus X gives us a chance to apply what we learn in class. We can use these concepts to build a working system that solves problems. This helps us see how the theory we study can be used to create something.}

\vspace{12mm}

\blueheading{State of the Art / Current solution (1 pt)}

\vspace{1mm}

\instruction{(Describe how the problem is solved today (if it is). Max 200 words).}

\answerbox[6cm]{Currently college-related information is managed through platforms instead of one unified system. WhatsApp groups are often used for announcements, discussions, notes and event updates. Instagram pages are used by clubs and societies to promote events and activities. Official information is usually shared through the college website, notice boards or separate communication channels. Students therefore have to check sources depending on what information they need.
\vspace{2mm}

Although these methods are helpful for purposes they can make campus information hard to manage and access. Important updates can get lost in group conversations. The same information might be shared times. Older information can stay mixed with updates. Some student-focused platforms offer features like communities, events, study groups and student connections.. They are not specifically built around the needs of our college. Campus X aims to bring all these activities in one organized platform. It will use DSA-based solutions, for searching, recommendations, ranking, notifications and information management.}

\vspace{12mm}
\newpage

\blueheading{Project Goals and Milestones (2 pts)}

\vspace{1mm}

\instruction{(Describe the project general goals. Include initial milestones as well any other milestones. Max 300 words).}

\answerbox[20cm]{Our primary goal is to develop \textbf{Campus X}, a student-focused platform for sharing information, connecting students, and accessing academic and campus-related resources. The project will demonstrate the practical use of Data Structures, Algorithms, and Object-Oriented Programming in C++. The main goals are:

\begin{itemize}
    \item Managing posts, announcements, events, clubs, and resources
    \item Providing friend and project partner recommendations
    \item Ranking content using popularity, relevance, importance, and recency
    \item Handling spam, duplicates, deleted posts, blocked users, and outdated information
\end{itemize}

Data Structures and Algorithms will form the core of the system:

\begin{itemize}
    \item \textbf{Hashing:} Duplicate detection and fast data access
    \item \textbf{Graph:} Friend and project partner recommendations
    \item \textbf{Priority Queue:} Emergency and important notifications
    \item \textbf{Trie:} Efficient searching
    \item \textbf{Queue and Stack:} Ordered and recent operations
    \item \textbf{Sorting and Searching:} Ranking and retrieving data
\end{itemize}

Object-Oriented Programming will be used to organize the system into manageable classes and modules.\\

\textbf{Planned Milestones:}

\begin{itemize}
    \item \textbf{Phase 1 -- Planning:} Finalize requirements, workflow, classes, and system design.
    \item Identify suitable Data Structures and Algorithms for each feature.
    \item \textbf{Phase 2 -- Development:} Implement the required DSA concepts and core Campus X modules in C++.
    \item Integrate posts, profiles, events, resources, recommendations, and notifications.
    \item \textbf{Phase 3 -- Testing:} Integrate the complete system and test all major features and edge cases.
    \item Optimize the algorithms and prepare the final demonstration.
\end{itemize}}

\vspace{12mm}
\newpage

\blueheading{Project Approach (3 pts)}

\vspace{1mm}

\instruction{(Describe how you plan to articulate and design a solution. Including platforms and technologies that you will use. Max}

\instruction{300 words).}

\answerbox[14cm]{\begin{enumerate}
    \item \textbf{Requirement Analysis and Planning} -- Find out what Campus X needs and set up how profiles, posts, announcements, events, clubs, resources, notifications, and recommendations will work together.

    \item \textbf{Data Structure Design} -- Represent the required data using structures like Hash Tables, Graphs, Tries, Queues, Stacks, and Priority Queues. Implement these structures using Object-Oriented Programming concepts in C++.

    \item \textbf{Post and Content Management} -- Build the post system with features such as creating, removing, sorting, and managing posts. Use hashing to check for duplicate posts and suitable structures to store and retrieve content.

    \item \textbf{Search and Recommendation Module} -- Implement searching using Tries and other searching techniques. Use Graphs to suggest friends and project partners based on connections and similar interests.

    \item \textbf{Ranking and Notification Module} -- Rank content based on popularity, relevance, importance, and recency. Use Priority Queues to ensure urgent and important notifications are shown first.

    \item \textbf{Exception and Edge-Case Handling} -- Handle real-life cases such as spam, repeated announcements, blocked users, deleted posts, outdated information, and restricted recommendations.

    \item \textbf{Module Integration and Testing} -- Combine all the modules into the Campus X application and test normal operations along with different edge cases.

    \item \textbf{Performance Analysis} -- Analyze the speed, memory usage, and efficiency of the implemented Data Structures and Algorithms and improve the system where required.
\end{enumerate}}

\newpage

% -------------------------------------------------
% PAGE 3
% -------------------------------------------------

\blueheading{System Architecture (High Level Diagram)(2 pts)}

\vspace{1mm}

\instruction{(Provide an overview of the system, identifying its main components and interfaces in the form of a diagram using a tool}

\instruction{of your choice).}

\answerbox[24.0cm]{
    \centering
    \includegraphics[height=24.0cm, keepaspectratio]{flowchart.png}
}

\vspace{12mm}

\blueheading{Project Outcome / Deliverables (1 pts)}

\vspace{1mm}

\instruction{(Describe what are the outcomes / deliverables of the project. Max 200 words).}

\answerbox[13cm]{The main goal of this project is to develop \textbf{Campus X}, a digital community platform designed specifically for college students.

\begin{itemize}[leftmargin=*,itemsep=2pt]
    \item It will provide a Campus Feed for announcements, achievements, discussions, opportunities and events.

    \item It will enable efficient searching of students, clubs, events, resources and campus information using Trie and Hashing.

    \item It will connect students through communities, relationships and the ``People You Should Know'' recommendation system using Graphs and BFS.

    \item It will manage resources and Lost \& Found information with efficient searching and duplicate detection.

    \item It  will manage events, notifications and emergency updates using Queues and Priority Queues.

    \item It will support Trending and Campus Leaderboard features using Heaps and Priority Queues.

    \item It  will organize departments, batches and clubs using Tree-based structures.

    \item It will provide Project Partner Matching based on student skills, connections and ranking algorithms.

    \item It  will handle real-world cases such as information, blocked users, deleted posts and changing student details.
\end{itemize}

The final deliverable will demonstrate how Data Structures and Algorithms can be applied to create a practical and well-organized digital platform for college communities.}

\vspace{12mm}

\blueheading{Assumptions}

\vspace{6mm}

\instruction{( Describe the assumptions ( if any ) you are making to solve the problem. Max 100 words )}

\answerbox[5cm]{
1. Only verified college students and authorized campus members will be able to use the Campus X platform.

2. Every student will have an account and a unique student identifier to manage their profile, connections, and activities.

3. Users must provide valid information when creating posts, events, resources, or other campus content.

4. Campus X will rely on Data Structures and Algorithms such as Tries, Hash Tables, Graphs, Heaps, Priority Queues, and Trees for efficient data management.

5. The system will preserve user and content history while handling blocked users, deleted posts, duplicate information, and changes in student details.}

\vspace{12mm}
\newpage

\blueheading{References}

\vspace{1mm}

\instruction{(Provide a list of resources or references you utilised for the completion of this deliverable. You may provide links).}

\answerbox[2.3cm]{
1. \textbf{GeeksforGeeks} -- for understanding Data Structures and Algorithms concepts.

2. \textbf{W3Schools} -- for SQL, JavaScript, and web development concepts.

3. \textbf{C++ Documentation} -- for C++ programming and STL concepts.

4. \textbf{Node.js \& Express.js Documentation} -- for backend development.

5. \textbf{MySQL Documentation} -- for SQL queries, database design, and data management.
}


\end{document}
